\chapter{The Language of Mathematics}\label{ch:language}

\begin{center}\itshape
The equation $x^2=1$ has a solution. It has a unique solution. Every real number solves it.\\
Three sentences, the same equation, three different questions.
\end{center}

Before you solve an equation, you need to know what its question means. Over the real numbers, $x^2=1$ has the two solutions $1$ and $-1$. Over the positive real numbers it has just one. The calculation has hardly changed. The sentence has. The domain, the words ``there is'' and ``every'', and the difference between existence and uniqueness all carry information. Learning to read that information is the first step toward writing a proof.

Figure~\ref{fig:lang-domains} draws the three readings. The rust marks are solutions. The cross is a real number that fails the equation.

\begin{figure}[ht]
\centering
\fig{ch00-domains}
\caption{The domain is part of the question. On all reals, $x^2=1$ has two solutions. On the positive reals, it has one. The test point $2$ shows that the equation is not solved by every real number.}
\label{fig:lang-domains}
\end{figure}

Here is a second question, in the same spirit. Someone claims that every increasing real function is positive. Increasing means that $x<y$ forces $f(x)\le f(y)$. Positive means that every value is greater than $0$. The examples offered are $x\mapsto x+1$ and $x\mapsto x+10$, for $x\ge 0$. How much evidence is enough? Neither example tests $x\mapsto x-1$ on the same domain. That function is increasing, and it takes the value $-1$. One counterexample defeats the claim. A thousand favourable examples would not prove it.

There is an old joke about a traveller who sees one sheep, and only one side of it, and announces that the sheep of the country are black. What was seen is smaller: at least one sheep, and at least one side of it black. The traveller said ``every'', and said it of the whole animal. Watch for that jump.

Trying examples helps you find a statement. A proof has to say why every case the words cover obeys them. The cube identity, for every $n$, is the first big sentence of that kind. This chapter fixes the words.

We start from arithmetic you already know. The symbols are there so you can see what was given, what you may choose, what you still have to show, and which choices may depend on which others.

A definition has to mean one thing, or a proof that uses the word is about something else. ``Every'' and ``some'' change the job even when the symbols look the same.

Leibniz hoped that a precise language would turn a dispute into a calculation. In \emph{De arte characteristica}, as printed in Gerhardt's edition of the philosophical writings, volume~7, page~200, the two disputants pick up their pens and say \emph{calculemus}: let us calculate. Clear writing helps. It does not pick the true sentence for you, and a line you cannot check is still a gap.

\portrait{leibniz}{Gottfried Wilhelm Leibniz (1646--1716). Portrait by Christoph Bernhard Francke, c.\,1700. Public domain, via Wikimedia Commons.}

\section{Definitions, assumptions and conclusions}

A definition and a theorem do different jobs. A definition gives a word an exact meaning. It is a test you can apply. It is not a claim waiting for a proof.

\begin{definitionx}[even and odd integers]
An integer $n$ is \emph{even} when $n=2k$ for some integer $k$. An integer $n$ is \emph{odd} when $n=2k+1$ for some integer $k$.
\end{definitionx}

To prove that $18$ is even, produce the integer $k=9$. To prove that $19$ is odd, produce the integer $k=9$, since $19=2\cdot 9+1$. To use the assumption that $n$ is even, write $n=2k$ for some integer $k$, and keep that integer in the argument. To use the assumption that $n$ is odd, write $n=2k+1$ for some integer $k$, and keep that integer in the argument.

Two different words can describe the same objects. That agreement then needs a proof of its own.

Names are cheap. Otto Blumenthal, in the life of Hilbert printed in Hilbert's \emph{Gesammelte Abhandlungen}, volume~3 (1935), pages~402--403, remembers a remark from a Berlin waiting room in 1891, after Hermann Wiener's lecture at Halle on the foundations of geometry. Hilbert's test, in that memory, was whether you could say ``tables, chairs, beer mugs'' wherever the text said ``points, lines, planes''. Try it on a definition. If the sentence stops meaning what you meant, it was using a picture you never wrote down. Blumenthal is remembering a conversation. The sentence does not appear in the \emph{Grundlagen der Geometrie} (1899).

An \emph{axiom} is an assumption that fixes the mathematical setting. A \emph{theorem} is a statement proved from the assumptions and from earlier results. The words \emph{proposition}, \emph{lemma} and \emph{corollary} also denote proved statements. A lemma serves another proof. A corollary follows readily from an earlier result. These names describe the role of the statement. They do not describe different standards of correctness. A \emph{conjecture} is a proposed statement for which a proof or a disproof is still sought. Many examples do not turn a conjecture into a theorem.

\begin{reminder}[our starting point]
We use ordinary classical logic, and the usual arithmetic and order of the natural numbers $\N=\{0,1,2,\dots\}$, the integers $\Z$, the rational numbers $\Q$ and the real numbers $\R$. The positive integers are $\N^*=\{1,2,\dots\}$. A rational number has the form $a/b$ with $a,b\in\Z$ and $b\ne 0$. The real-number axioms include completeness. Its consequences are developed in Chapter~\ref{ch:sequences}. Well ordering of $\N$ is stated and used in Chapter~\ref{ch:sums}: every nonempty subset of $\N$ has a least element. We use sets and functions in their ordinary mathematical sense. We do not construct set theory from formal axioms. Later chapters state the elementary-function and integration facts they assume. Nothing here requires a previous course in proofs.
\end{reminder}

In a theorem of the form ``if $H$, then $C$'', the statement $H$ is the \emph{hypothesis} and $C$ is the \emph{conclusion}. A proof may use $H$, the definitions, the stated foundations and results already established. It may not use $C$ as an extra assumption and then announce that $C$ has been proved. Working backwards from $C$ can be useful on scrap paper. The final argument has to justify every step from the information already available.

\begin{examplex}[unpacking a definition]
The square of every even integer is even. Let $n$ be an even integer. By definition, $n=2k$ for an integer $k$. Then
\[
n^2=(2k)^2=4k^2=2\cdot(2k^2).
\]
The number $2k^2$ is an integer, because $k$ is an integer. The definition now says that $n^2$ is even. The witness for the hypothesis was $k$. The witness for the conclusion is $2k^2$. The proof is a calculation guided by two uses of one definition.
\end{examplex}

You unpacked a word, calculated, and packed the word up again. That is a direct proof. The sentence that held the steps together was an ``if'': if $n$ is even, then $n^2$ is even. What ``if'' promises, and what ``and'' and ``or'' promise, still has to be said.

\section{Statements, connectives and implication}

A \emph{statement} is a sentence with a definite truth value in its stated setting. The sentence ``$7>3$'' is true. The sentence ``$7<3$'' is false. The expression ``$x>3$'' becomes a statement after $x$ is specified, or after a quantifier binds it. Until then we call it a \emph{predicate} depending on $x$. Fix the domain before you treat an expression as a predicate.

For statements $P$ and $Q$, the connective $P\land Q$ means ``$P$ and $Q$''. It is true exactly when both are true. The connective $P\lor Q$ means ``$P$ or $Q$''. It is true when at least one is true, and it remains true when both are true. This is the inclusive mathematical meaning of ``or''. A menu that says ``soup or salad'' is often read as a choice of one plate. Mathematical ``or'' lets you take both. The four cases are:
\[
\begin{array}{c|c|c}
P&Q&P\lor Q\\\hline
\text{true}&\text{true}&\text{true}\\
\text{true}&\text{false}&\text{true}\\
\text{false}&\text{true}&\text{true}\\
\text{false}&\text{false}&\text{false}
\end{array}
\]
The top row is the one a menu often hides: both parts true, and the ``or'' still true. Exercise~\ref{ex:lang-vacuous} asks you to use that reading. The negation $\neg P$ means ``$P$ is false''. In prose we usually write these words rather than the symbols.

The implication $P\Rightarrow Q$ says that $P$ cannot be true while $Q$ is false. It has the same truth value as $(\neg P)\lor Q$. The four cases are worth seeing once:
\[
\begin{array}{c|c|c}
P&Q&P\Rightarrow Q\\\hline
\text{true}&\text{true}&\text{true}\\
\text{true}&\text{false}&\text{false}\\
\text{false}&\text{true}&\text{true}\\
\text{false}&\text{false}&\text{true}
\end{array}
\]
An implication promises something when its hypothesis holds. It promises nothing about the cases where that hypothesis fails. Thus ``if $x>2$, then $x^2>4$'' is true for every real $x$, including $x=0$. At $x=0$ the hypothesis is false, so that point does not provide a counterexample. This convention is called \emph{vacuous truth}.

Every elephant in this room is an even integer. If the room has no elephant, there is no elephant that fails the test, so the sentence is true. The empty set will be a subset of every set for the same reason. Older logic often read ``every $A$ is $B$'' as already saying that some $A$ exists. Here a universal sentence fails only when you can point at a counterexample.

An implication does not assert a cause. It also does not let you reverse the arrow. From $x=1\Rightarrow x^2=1$, the converse $x^2=1\Rightarrow x=1$ does not follow. The number $x=-1$ defeats the converse: $(-1)^2=1$, and $-1\ne 1$. A false hypothesis does not force a false conclusion. From ``if $n$ is a multiple of $4$, then $n$ is even'', we cannot deduce that an integer which is not a multiple of $4$ fails to be even. The integer $6$ is even, and $6$ is not a multiple of $4$. Here ``$a$ is a multiple of $b$'' means $a=bk$ for an integer $k$.

The equivalence $P\Leftrightarrow Q$ means that both $P\Rightarrow Q$ and $Q\Rightarrow P$ hold. It is read ``$P$ if and only if $Q$'', often shortened to ``iff''. For a real number $x$, the equivalence $x^2=0\Leftrightarrow x=0$ holds. One direction substitutes $x=0$ and gets $0^2=0$. In the other direction, suppose $x^2=0$ and $x\ne 0$. The real numbers allow division by the nonzero number $x$, and dividing $x^2=0$ by $x$ gives $x=0$. That contradicts $x\ne 0$. So the extra supposition $x\ne 0$ is impossible, and $x=0$.

If $P\Rightarrow Q$, then $P$ is a \emph{sufficient} condition for $Q$, and $Q$ is a \emph{necessary} condition for $P$. Being a multiple of $4$ is sufficient for being even. Being even is necessary for being a multiple of $4$. To claim a necessary and sufficient condition is to claim an equivalence, so both directions have to be justified.

\begin{think}
For a real number $x$, is $x>1$ necessary for $x^2>1$? Is it sufficient?
Test a negative number before you decide. Which necessary and sufficient condition for $x^2>1$ uses both positive and negative numbers?
\end{think}

Here is the check. First, sufficiency. Suppose $x>1$. Then $x-1>0$. Also $x+1>1+1=2>0$. The product of two positive numbers is positive, so
\[
x^2-1=(x-1)(x+1)>0,
\]
and therefore $x^2>1$. So $x>1$ is sufficient for $x^2>1$.

It is not necessary. At $x=-2$ we have $(-2)^2=4>1$, and $-2>1$ is false.

The same check works for $x<-1$. Suppose $x<-1$. Then $x+1<0$. Also $x-1<x+1<0$. Both factors are negative, and the product of two negative numbers is positive, so
\[
x^2-1=(x-1)(x+1)>0,
\]
and therefore $x^2>1$. So $x<-1$ is sufficient as well.

For the converse, suppose $-1\le x\le 1$. If $0\le x\le 1$, then $x-1\le 0$ and $x+1>0$, so $(x-1)(x+1)\le 0$. If $-1\le x<0$, then $x-1<0$ and $x+1\ge 0$, so the product is again at most $0$. In both cases $x^2\le 1$. Thus $x^2>1$ holds exactly when $x>1$ or $x<-1$. So $x>1$ or $x<-1$ is necessary and sufficient for $x^2>1$. Figure~\ref{fig:lang-rays} draws the two rays on which the inequality holds, and the segment on which it fails. Exercise~\ref{ex:lang-implications} asks you to carry out the twin argument for $x^2>4$.

\begin{figure}[ht]
\centering
\fig{ch00-rays}
\caption{The inequality $x^2>1$ holds on the two rust rays, where $x<-1$ and where $x>1$. The right-hand ray is the condition $x>1$: sufficient, and not necessary. On the charcoal segment, including the endpoints $\pm 1$, one has $x^2\le 1$.}
\label{fig:lang-rays}
\end{figure}

The connecting words say how statements combine. They still do not say who is allowed to choose a number. That comes next.

\section{Quantifiers: who chooses what?}

The \emph{universal quantifier} $\forall$ means ``for every''. The \emph{existential quantifier} $\exists$ means ``there exists''. The signs are younger than the words. The sign $\exists$ appears in Peano's \emph{Formulaire de math\'ematiques} of 1897, and Russell used it to bind a variable. Gerhard Gentzen introduced $\forall$ in print in 1935, in \emph{Untersuchungen \"uber das logische Schlie\ss en} (\emph{Mathematische Zeitschrift} 39), as the mate of $\exists$, which he records taking from Russell. Frege's \emph{Begriffsschrift} (1879) had a notation for generality, a way to write ``for every''. ``There exists'' was written as ``not every $\dots$ not''. There was no separate sign for existence. Treat $\forall$ and $\exists$ as shorthand, and read them in English. Keep the domain in the sentence:
\[
\forall x\in\R,\quad x^2\ge 0;
\qquad
\exists x\in\R,\quad x^2=2.
\]
The first asks for a proof that covers an arbitrary real $x$. The second asks for at least one witness. A real witness for $x^2=2$ exists by the intermediate value theorem, Theorem~\ref{thm:ivt} in Chapter~\ref{ch:sequences}, applied to the continuous function $x\mapsto x^2-2$ on the interval $[1,2]$. This chapter does not construct that witness. Asking for a rational witness is a different question. A real solution does not settle it, and this chapter does not settle it either.

For an example you can finish now, compare domains. The statement $\exists x\in\R,\ 2x=1$ is true, with witness $x=1/2$. The statement $\exists x\in\Z,\ 2x=1$ is false. If $k\ge 1$ is an integer, then $2k\ge 2$. If $k\le 0$ is an integer, then $2k\le 0$. So twice an integer never lies strictly between $0$ and $2$, and it cannot equal $1$.

To see how a witness may depend on earlier choices, compare
\[
\forall x\in\R\ \exists y\in\R:\ x+y=0
\quad\text{and}\quad
\exists y\in\R\ \forall x\in\R:\ x+y=0.
\]
The first is true. After $x$ is given, choose $y=-x$. This $y$ is allowed to depend on $x$. Check: $x+(-x)=0$. The second is false. A single $y$ would have to satisfy both $0+y=0$ and $1+y=0$. The first equation forces $y=0$. The second forces $y=-1$. Those demands disagree. The order of the quantifiers says which information is available when a choice is made.

Figure~\ref{fig:lang-quantifiers} draws the two sentences. On the left, each $x$ receives its own $y$. On the right, one $y$ is chosen first and then has to work for every $x$.

\begin{figure}[ht]
\centering
\fig{ch00-quantifiers}
\caption{Two orders, two meanings. On the left, $y$ is chosen after $x$, so $y=-x$ solves $x+y=0$ for every real $x$. On the right, $y$ is chosen first. It cannot be both $0$ and $-1$.}
\label{fig:lang-quantifiers}
\end{figure}

Adjacent universal quantifiers may be exchanged. ``For every $x$ and every $y$'' asks for all pairs, in either order. Adjacent existential quantifiers may also be exchanged. A successful pair supplies both witnesses. Mixed quantifiers generally cannot be exchanged. The pair of sentences above is the example. When you read a long statement, mark the variables that are chosen universally, and draw an arrow from each existential choice back to the earlier variables it may use.

Later, the order $\forall\varepsilon>0\ \exists N\ \forall n\ge N$ will make convergence precise. The number $N$ may depend on $\varepsilon$. It may not depend on the later index $n$. You do not need that definition yet. You need the habit of reading the order.

The \emph{scope} of a quantifier is the part of the sentence to which it applies. Use punctuation or parentheses when the scope could be read in two ways. The statement $\forall x\in\R,\ (x>0\Rightarrow x^2>0)$ quantifies an implication. It does not say that every real number is positive. Quantifiers over an empty domain follow the same convention as implication. ``Every element of the empty set has property $P$'' is true, because there is no element at which $P$ fails. ``Some element of the empty set has property $P$'' is false, because there is no witness.

\subsection*{Negation without changing the question}

Negation reverses each quantifier and negates the predicate. It keeps the order of the variables.

\begin{proposition}[negating a quantified statement]\label{prop:lang-negation}
For a set $E$ and a predicate $P$,
\begin{align*}
\neg(\forall x\in E,\ P(x))&\ \Longleftrightarrow\ \exists x\in E,\ \neg P(x),\\
\neg(\exists x\in E,\ P(x))&\ \Longleftrightarrow\ \forall x\in E,\ \neg P(x).
\end{align*}
The same meanings give the rules for connectives:
\begin{align*}
\neg(P\land Q)&\ \Longleftrightarrow\ (\neg P)\lor(\neg Q),\\
\neg(P\lor Q)&\ \Longleftrightarrow\ (\neg P)\land(\neg Q),\\
\neg(P\Rightarrow Q)&\ \Longleftrightarrow\ P\land(\neg Q).
\end{align*}
\end{proposition}
\begin{proof}
The universal claim fails exactly when some element of $E$ fails $P$. That element is a witness for the existential statement on the first line. The existential claim fails exactly when no element of $E$ satisfies $P$, which means that every element fails $P$.

If $E$ is empty, there is no witness. Every existential statement about $E$ is then false, and every universal statement about $E$ is true. Both lines hold in that case as well.

For the connectives, $P\land Q$ fails when at least one of $P,Q$ fails, and $P\lor Q$ fails only when both fail. The implication $P\Rightarrow Q$ fails only in the row of the table where $P$ is true and $Q$ is false. That row is exactly $P\land(\neg Q)$.
\end{proof}

For example, ``$x$ does not satisfy $0<x<1$'' means $x\le 0$ or $x\ge 1$. The negation of an ``and'' is an ``or''. The endpoints belong to the negation: the original strict inequalities fail at $0$ and at $1$.

Proposition~\ref{prop:lang-negation} has five rules: two for the quantifiers, and three for ``and'', ``or'' and ``if''. Two of them are De Morgan's laws for statements. The negation of an ``and'' is an ``or'' of negations, and the negation of an ``or'' is an ``and'' of negations. The same pattern for sets comes later, in Proposition~\ref{prop:lang-sets}: the complement of a union is the intersection of the complements, and the complement of an intersection is the union of the complements. In \emph{Formal Logic} (1847) Augustus De Morgan said it for names: when you pass to the contrary, ``and'' and ``or'' change places. A form of that swap is older, in medieval logic. He was trying to calculate with logical words.

George Boole's \emph{The Mathematical Analysis of Logic} (1847) and \emph{An Investigation of the Laws of Thought} (1854) treat classes as things you can calculate with. Watch the plus sign. Boole added two classes only when they did not overlap. Our ``or'' is the inclusive one in the table for $P\lor Q$: the row where both parts are true still counts as true. Calculate, then check the line against what the words mean.

\portrait{boole}{George Boole (1815--1864). Engraving in \emph{The Illustrated London News}, 21 January 1865. Public domain, via Wikimedia Commons.}

Read a negation from the outside inward. The sentence ``for every real $x$ there is a real $y$ with $y>x$'' is true: given $x$, take $y=x+1$. Its negation is
\[
\exists x\in\R\ \forall y\in\R,\quad y\le x.
\]
That sentence says there is a greatest real number. It is false, because $y=x+1$ is always a larger real. Notice what the negation does. It changes $\forall$ to $\exists$ and $\exists$ to $\forall$, and it replaces $y>x$ by $y\le x$. It does not insert a minus sign, and it does not keep the outer quantifier as ``for every''.

\begin{examplex}[a negation that looks plausible and is wrong]
The negation of ``for every real $x$ there is a real $y$ with $y>x$'' is not ``for every real $x$ there is a real $y$ with $y\le x$''. That second sentence is true: given $x$, take $y=x$. The original sentence is also true. A negation of a true sentence has to be false. The correct negation is the sentence displayed above, which asks for one $x$ that beats every $y$. Write the negation before you decide whether it holds.
\end{examplex}

You can say ``every'' and ``some'', and you can deny them. Denial keeps the order of the variables and flips each quantifier as it goes. A proof has to be just as literal: the sentence tells you whether you may choose a witness or must take an arbitrary element.

\section{Choosing a proof method}

The shape of the sentence often tells you the first line. To prove ``for every $x\in E$, $P(x)$'', let $x$ be an arbitrary element of $E$, and establish $P(x)$ without adding a special restriction on $x$. To prove ``there exists'', construct a witness, or prove that a witness must exist. To prove an equivalence, write both directions. Equality of two sets, defined in the next section, has the same shape: check membership each way. The words tell you that first line. A clever step, if you need one, comes after it.

\subsection*{Direct proof and proof by cases}

A \emph{direct proof} starts from the hypotheses and reaches the conclusion. The even-square argument was direct. Here is another. If $a<b$ are real numbers, there is a real number strictly between them.

Set $m=(a+b)/2$. Then
\[
m-a=\frac{b-a}{2}.
\]
The numerator $b-a$ is positive because $a<b$, and the denominator $2$ is positive, so the quotient is positive. Thus $m>a$. The same computation gives $b-m=(b-a)/2>0$, so $m<b$. Therefore $a<m<b$. The proof supplies a witness and checks the two inequalities the word ``between'' asks for.

When the hypotheses offer alternatives, the proof can handle them separately. For every real $x$, the inequality $|x|\ge x$ holds. If $x\ge 0$, then $|x|=x$, so equality holds. If $x<0$, then $|x|=-x$, and $-x>0>x$, so $|x|>x$. A proof by cases has to cover every allowed case. Splitting into $x>0$ and $x<0$ would omit $x=0$. The split $x\ge 0$ and $x<0$ covers every real.

\subsection*{Contrapositive}

The implication $P\Rightarrow Q$ is equivalent to its \emph{contrapositive} $\neg Q\Rightarrow\neg P$. Both fail in exactly one case: $P$ true and $Q$ false. The contrapositive offers a different starting point. It does not change the theorem.

For example, if $x$ is irrational, then $x+1$ is irrational. An irrational number is a real number that is not rational. We prove the contrapositive: if $x+1$ is rational, then $x$ is rational. Suppose $x+1=a/b$ with $a,b\in\Z$ and $b\ne 0$. Then
\[
x=(x+1)-1=\frac{a}{b}-\frac{b}{b}=\frac{a-b}{b}.
\]
The numerator $a-b$ is an integer, and the denominator $b$ is a nonzero integer, so $x$ is rational. The assumption we used is the negation of the desired conclusion. The conclusion we reached is the negation of the original hypothesis.

The converse is a different sentence: ``if $x+1$ is irrational, then $x$ is irrational''. That converse happens to be true as well, by a similar subtraction. Truth of the converse is not what the contrapositive proof established. The contrapositive of $P\Rightarrow Q$ is $\neg Q\Rightarrow\neg P$. The converse is $Q\Rightarrow P$.

\subsection*{Contradiction}

A \emph{proof by contradiction} assumes the hypotheses together with the negation of the desired conclusion, and derives an impossibility. Classical logic then rules out that negation.

There is no smallest positive real number. Suppose, for a contradiction, that $a>0$ is a smallest positive real. The number $a/2$ is positive, because $a>0$ and $2>0$. It is strictly smaller than $a$, because
\[
a-\frac{a}{2}=\frac{a}{2}>0.
\]
So $a/2$ is a positive real smaller than $a$. This contradicts the choice of $a$. The two statements cannot both hold: $a$ is no larger than any positive real, and $a/2$ is a positive real below $a$.

Here is what a contradiction costs, in ordinary arithmetic. Suppose $0=1$. Take any integer $n$ and multiply both sides by $n$. Then $0=n$. Every integer equals $0$. One false equation, and the arithmetic gives you all the others. Classical logic does the same with sentences: from a statement and its negation, every statement follows. Use this when you need to throw out a bad assumption.

A contrapositive proof of $P\Rightarrow Q$ proves the new implication $\neg Q\Rightarrow\neg P$. A contradiction proof may begin with $P$ and $\neg Q$ together and derive any impossibility. Use whichever one gives you something you can calculate with. If the calculation was already direct, adding ``suppose not'' does not make it clearer.

\subsection*{Counterexamples, existence and uniqueness}

To disprove $\forall x\in E,\ P(x)$, give one $x\in E$ and show that $P(x)$ fails. For an implication, the example must satisfy the hypothesis. To disprove ``every positive real $x$ has $x^2>x$'', the number $x=-1$ is irrelevant, because $-1$ is not positive. The number $x=1/2$ works: it is positive, and $(1/2)^2=1/4$, which is strictly less than $1/2$. Counterexamples often live at $0$, at an endpoint, at equal parameters, or in a small structure where an informal picture breaks.

An existence and uniqueness claim has two parts. The notation $\exists!\, x\in E,\ P(x)$ abbreviates
\[
\exists x\in E:\quad P(x)\ \text{and}\ \forall y\in E,\ (P(y)\Rightarrow y=x).
\]
Finding a candidate proves neither part by itself. For the equation $3x+2=8$ over $\R$, rearrangement suggests $x=2$. Substitution checks existence: $3\cdot 2+2=8$. For uniqueness, suppose $3y+2=8$. Subtract $2$ to get $3y=6$. The number $3$ is nonzero, so it has a multiplicative inverse in $\R$. Multiplying by that inverse gives $y=2$. In a harder problem, an existence proof may identify no explicit formula. The intermediate value theorem will provide an important example. Whenever the book claims uniqueness, look for the step that excludes a second solution.

\begin{reminder}[errors that survive a convincing calculation]
An argument can fail by proving the converse, by choosing an existential witness before the data it must depend on, by checking only examples, by assuming the conclusion, or by using an expression outside its domain. Dividing by a quantity requires knowing that the quantity is nonzero. Squaring an equation can introduce extra solutions. Once non-negative square roots exist, the symbol $\sqrt{\ }\,$ returns that root, so $\sqrt{x^2}=|x|$. It does not return both solutions of $t^2=x^2$. Finally, write each ``therefore'' as a small question: does the information on the page force the next assertion?
\end{reminder}

Each of those shapes already talked about objects: an integer $n$, a real $x$, a witness you had to produce. From the quantifiers onward, $E$ was the collection those objects belong to. That collection is a set, and membership is the test for whether your witness is in it. The workshop later asks the image question with two subsets of one set.


\section{Sets: making the domain visible}

A \emph{set} is a collection of distinct objects, called its \emph{elements}. We write $x\in A$ for membership and $x\notin A$ for non-membership. Listing braces ignore order and repetition, so $\{1,2,2\}=\{2,1\}$. The empty set $\varnothing$ has no elements. Set-builder notation records a condition. For example, $\{x\in\R:x^2=1\}=\{-1,1\}$. The colon means ``such that''.

We write $A\subseteq B$ when every element of $A$ belongs to $B$:
\[
A\subseteq B\quad\Longleftrightarrow\quad
\forall x,\ (x\in A\Rightarrow x\in B).
\]
Equality of sets means equality of membership. So $A=B$ is equivalent to the two inclusions $A\subseteq B$ and $B\subseteq A$. To prove two sets equal, prove both inclusions. A \emph{proper} subset, written $A\subsetneq B$, is a subset unequal to $B$. Some authors use $\subset$ for either meaning. In this book, $\subseteq$ allows equality.

For every set $A$, the inclusion $\varnothing\subseteq A$ holds. The implication $x\in\varnothing\Rightarrow x\in A$ has no counterexample, because nothing lies in $\varnothing$. This is vacuous truth from the previous section, applied to membership.

The \emph{union}, \emph{intersection} and \emph{difference} are
\begin{align*}
A\cup B&=\{x:x\in A\ \text{or}\ x\in B\},\\
A\cap B&=\{x:x\in A\ \text{and}\ x\in B\},\\
A\setminus B&=\{x:x\in A\ \text{and}\ x\notin B\}.
\end{align*}
These descriptions take their elements from the sets already under discussion. For a complement we fix an \emph{ambient set} $U$ containing $A$, and we write $A^c=U\setminus A$. The complement of $\{1,2\}$ in $\{1,2,3\}$ is $\{3\}$. The complement of $\{1,2\}$ in $\N$ is a different set. We do not assume a set of all mathematical objects.

On 16 June 1902 Bertrand Russell wrote to Gottlob Frege that there is no class, taken as a totality, of the classes which do not belong to themselves. If $w$ were that class, then each of $w\in w$ and $w\notin w$ would force the other. Frege got the letter while the second volume of the \emph{Grundgesetze der Arithmetik} was in press. The letter is printed in van Heijenoort's \emph{From Frege to G\"odel}. Repairing that belongs to a later set theory, and we do not build it here. What we do instead is smaller. Every complement is taken inside a set you have named, and a condition is a set only when you have said which objects it talks about. The last challenge of the chapter makes the same point with power sets: no function from a set to its power set hits every subset.

For a family of sets $(A_i)_{i\in I}$, membership in $\bigcup_{i\in I}A_i$ means membership in at least one $A_i$. Membership in $\bigcap_{i\in I}A_i$ means membership in every $A_i$. When the intersection is taken inside a fixed ambient set $U$, the empty union is $\varnothing$ and the empty intersection is $U$. These are conventions. They match existential and universal quantification over an empty index set: an empty ``there exists'' is false, and an empty ``for every'' is true.

\begin{proposition}[set operations and logic]\label{prop:lang-sets}
For subsets $A,B,C$ of a fixed set $U$,
\begin{align*}
A\cap(B\cup C)&=(A\cap B)\cup(A\cap C),\\
(A\cup B)^c&=A^c\cap B^c,\qquad
(A\cap B)^c=A^c\cup B^c.
\end{align*}
\end{proposition}
\begin{proof}
Take $x\in U$. We compare membership, one identity at a time.

\emph{Distributive law.} The element $x$ lies in $A\cap(B\cup C)$ exactly when $x\in A$ and at least one of $x\in B$, $x\in C$ holds. If $x\in B$, then $x\in A\cap B$. If $x\in C$, then $x\in A\cap C$. In either case $x$ lies in $(A\cap B)\cup(A\cap C)$. Conversely, if $x$ lies in that union, then $x$ lies in $A\cap B$ or in $A\cap C$. In the first sub-case $x\in A$ and $x\in B$. In the second, $x\in A$ and $x\in C$. Either way, $x\in A$ and $x\in B\cup C$.

\emph{Complement of a union.} The element $x$ lies outside $A\cup B$ exactly when it lies outside $A$ and outside $B$. That is membership in $A^c\cap B^c$.

\emph{Complement of an intersection.} The element $x$ fails to lie in both $A$ and $B$ exactly when at least one of those memberships fails. That is membership in $A^c\cup B^c$.

The two complement identities are the ``and'' and ``or'' rules from Proposition~\ref{prop:lang-negation}, read as membership. The distributive law is a different pattern: ``and'' passing across ``or''.
\end{proof}

A small check: let $U=\{1,2,3\}$, $A=\{1,2\}$, $B=\{2\}$ and $C=\{3\}$. Then $B\cup C=\{2,3\}$ and $A\cap(B\cup C)=\{2\}$. Also $A\cap B=\{2\}$ and $A\cap C=\varnothing$, so the union of those intersections is $\{2\}$. The two sides agree. Also $A\cup B=\{1,2\}$, so its complement in $U$ is $\{3\}$. And $A^c=\{3\}$, $B^c=\{1,3\}$, so $A^c\cap B^c=\{3\}$. The two sides agree again.

Figure~\ref{fig:lang-demorgan} draws the complement law with the numbers left out. The rust region is the set of points in neither circle.

\begin{figure}[ht]
\centering
\fig{ch00-demorgan}
\caption{The rust region lies in neither $A$ nor $B$. A point there lies in $(A\cup B)^c$, and it lies in $A^c\cap B^c$. Proposition~\ref{prop:lang-sets} says these are the same set of points of $U$.}
\label{fig:lang-demorgan}
\end{figure}

You will use this move often. Draw the condition for $x$, then rewrite that condition. Do not shift the symbols first and hunt for a meaning afterwards.

\subsection*{Power sets and ordered pairs}

The \emph{power set} $\mathcal P(A)$ is the set of all subsets of $A$. For example,
\[
\mathcal P(\{a,b\})=\{\varnothing,\{a\},\{b\},\{a,b\}\}.
\]
There are four subsets: include $a$ or not, and include $b$ or not. Membership and inclusion now live side by side. The statement $\{a\}\in\mathcal P(A)$ means the same thing as $\{a\}\subseteq A$. The set $\varnothing$ differs from $\{\varnothing\}$. The first has no elements. The second has one element, namely $\varnothing$.

An \emph{ordered pair} $(a,b)$ records a first coordinate and a second coordinate. The equality test is
\[
(a,b)=(c,d)\quad\text{exactly when}\quad a=c\ \text{and}\ b=d.
\]
A two-element set forgets order and cannot list a repeated coordinate twice. An ordered pair does both. The \emph{Cartesian product} is
\[
A\times B=\{(a,b):a\in A,\ b\in B\}.
\]
Thus $\{0,1\}\times\{2\}=\{(0,2),(1,2)\}$, while $\{2\}\times\{0,1\}=\{(2,0),(2,1)\}$. These products are different sets. Repeated products give tuples $(x_1,\dots,x_n)$. In particular, $\R^2$ is the coordinate plane, and $\R^n$ is the set of real $n$-tuples. Products will describe successive choices in counting, and successive trials in probability.

Sets name the objects. A function is a rule that sends each object of one set to an object of another. That comes next.

\section{Functions: the rule and its two sets}

\begin{definitionx}[function]
A \emph{function} $f:E\to F$ assigns to every element $x$ of its \emph{domain} $E$ exactly one element $f(x)$ of its \emph{codomain} $F$. The symbol $x\mapsto f(x)$ describes the assignment. The letter $f$ names the function. The expression $f(x)$ names one value.
\end{definitionx}

One can identify $f$ with its \emph{graph}, the subset $\{(x,f(x)):x\in E\}$ of $E\times F$, together with the specified domain and codomain. The graph meets each vertical line $\{x\}\times F$ in exactly one point. That is the same demand as ``exactly one value''.

The word \emph{function} is older than the box. In 1673 Leibniz used it for a geometric quantity read off a curve. In the \emph{Introductio in analysin infinitorum} (1748), Euler said a function of a variable was an analytic expression built from that variable and from constants. In the \emph{Institutiones calculi differentialis} (1755) he widened it: a quantity fixed once other quantities are given. The box is the later meaning, the one Dirichlet was using in 1837. Each allowed input gets one output. A single formula is not required. The rule is the definition. Use a formula when you have one.

\begin{center}
\portraitbox{0.48\linewidth}{euler}{Leonhard Euler (1707--1783). Portrait by Jakob Emanuel Handmann, 1753, Kunstmuseum Basel. Public domain, via Wikimedia Commons.}
\end{center}

A formula alone need not specify a function. The formula $1/x$ cannot define a function $\R\to\R$, because it has no value at $0$. It does define a function $\R\setminus\{0\}\to\R$. The notation ``$f(x)=\pm\sqrt{x}$'' does not select one value for a positive $x$. The formula $x^2$ does define a function on all of $\R$, and it also defines a function on the smaller domain $[0,\infty)=\{x\in\R:x\ge 0\}$.

\begin{reminder}[interval notation]
The interval $[a,b]$ includes both endpoints. The interval $(a,b)$ includes neither. A square bracket at one end includes only that endpoint, so $[a,b)$ includes $a$ and excludes $b$. An infinite end is always excluded. Later chapters also write $]a,b[$ for $(a,b)$, $[a,b[$ for $[a,b)$, and $]a,b]$ for $(a,b]$. An outward-facing square bracket marks an excluded endpoint.
\end{reminder}

For $A\subseteq E$, the \emph{image} of $A$ is
\[
f(A)=\{f(x):x\in A\}=\{y\in F:\exists x\in A,\ f(x)=y\}.
\]
The image of the whole domain, $f(E)$, is also called the \emph{range}. It may be smaller than the codomain. For $B\subseteq F$, the \emph{inverse image}, or \emph{preimage}, is
\[
f^{-1}(B)=\{x\in E:f(x)\in B\}.
\]
This notation makes sense for every function. It does not claim that $f$ has an inverse function. For $f:\R\to\R$ given by $f(x)=x^2$,
\[
f(\{-2,1\})=\{4,1\},\qquad f^{-1}(\{1\})=\{-1,1\},\qquad f^{-1}((-\infty,0))=\varnothing.
\]
The last set is empty because a square is never negative.

\begin{proposition}[inverse images preserve set operations]\label{prop:lang-preimages}
If $f:E\to F$ and $B,C\subseteq F$, then
\begin{align*}
f^{-1}(B\cup C)&=f^{-1}(B)\cup f^{-1}(C),\\
f^{-1}(B\cap C)&=f^{-1}(B)\cap f^{-1}(C),\\
f^{-1}(F\setminus B)&=E\setminus f^{-1}(B).
\end{align*}
\end{proposition}
\begin{proof}
Take $x\in E$. Membership in $f^{-1}(B\cup C)$ means $f(x)\in B\cup C$, which means $f(x)\in B$ or $f(x)\in C$. That is membership in $f^{-1}(B)\cup f^{-1}(C)$.

Replace ``or'' by ``and''. Membership in $f^{-1}(B\cap C)$ means $f(x)\in B$ and $f(x)\in C$, which is membership in the intersection of the preimages.

For the third identity, $f(x)\notin B$ means that $x$ fails to lie in $f^{-1}(B)$. The complement on the left is taken in $F$. The complement on the right is taken in $E$.
\end{proof}

A check with $f(x)=x^2$ on $\R$: let $B=[0,\infty)$ and $C=(-\infty,0]$. Then $B\cap C=\{0\}$, and $f^{-1}(\{0\})=\{0\}$. Also $f^{-1}(B)=\R$ and $f^{-1}(C)=\{0\}$, so their intersection is $\{0\}$. The two sides agree.

These equalities are why an event described by a numerical measurement will be an inverse image. ``The measured value lies in $B$'' is a condition on the original outcomes. An image collects values, and it can forget which inputs produced them. That is why images do not always preserve intersections. The workshop at the end of the chapter makes the failure precise. Figure~\ref{fig:language-fibres} shows the two ways a function can be less neat than a pairing: two inputs can share a value, and a codomain element can be missed.

\begin{figure}[ht]
\centering
\fig{ch00-fibres}
\caption{A function can send two inputs to one value and miss another value. Here $f(x_1)=f(x_2)=a$ and $f(x_3)=b$, while $c$ has no preimage. The dashed set is one fibre. The image $\{a,b\}$ is smaller than the codomain $F=\{a,b,c\}$.}
\label{fig:language-fibres}
\end{figure}

The fibre in that figure is the preimage of a single value. The next section turns ``two inputs, one value'' and ``a missed value'' into tests with names.

\section{Injectivity, surjectivity and inverse functions}

\begin{definitionx}[injective, surjective, bijective]
A function $f:E\to F$ is \emph{injective} when equal outputs force equal inputs:
\[
\forall x,x'\in E,\quad f(x)=f(x')\Rightarrow x=x'.
\]
It is \emph{surjective} when every element of the specified codomain occurs as a value:
\[
\forall y\in F\ \exists x\in E,\quad f(x)=y.
\]
It is \emph{bijective} when it is both injective and surjective.
\end{definitionx}

A bijection pairs each input with one output and each output with one input. Counting will use such pairings to show that two finite sets have the same number of elements.

The function $x\mapsto x^2$ from $\R$ to $\R$ has neither property. The inputs $1$ and $-1$ have the same output, so the map is not injective. The number $-1$ has no preimage, so the map is not surjective onto $\R$.

Restrict the domain to $[0,\infty)$. The map becomes injective. Here is the check, in steps. Suppose $x,x'\ge 0$ and $x^2=(x')^2$. Then
\[
x^2-(x')^2=(x-x')(x+x')=0.
\]
So $x-x'=0$ or $x+x'=0$. If $x-x'=0$, then $x=x'$. If $x+x'=0$, then $x'=-x$. Both numbers are at least $0$, so $-x=x'\ge 0$, which means $x\le 0$. Combined with $x\ge 0$ this forces $x=0$, and then $x'=0$. In both cases the inputs agree.

Surjectivity onto $[0,\infty)$ is a different claim. It says that every $y\ge 0$ has a non-negative square root. That existence is an application of the intermediate value theorem, Theorem~\ref{thm:ivt} in Chapter~\ref{ch:sequences}, after that chapter records that a polynomial is continuous. For $y>0$ the function $t\mapsto t^2$ runs from $0$ to $(y+1)^2>y$ on $[0,y+1]$. Until that theorem, do not treat $x\mapsto x^2$ as a bijection of $[0,\infty)$ onto $[0,\infty)$. The example separates the job of the domain from the job of the codomain.

Given $f:E\to F$ and $g:F\to G$, their \emph{composition} is the function $g\circ f:E\to G$ defined by $(g\circ f)(x)=g(f(x))$. Read the right-hand function first: $f$ acts, then $g$. For instance, if $f(x)=x+1$ and $g(x)=x^2$ on $\R$, then $g\circ f$ sends $x$ to $(x+1)^2$, while $f\circ g$ sends $x$ to $x^2+1$. At $x=1$ the first value is $4$ and the second is $2$. Composition is associative. For a third function $h$,
\begin{align*}
\bigl(h\circ(g\circ f)\bigr)(x)
&=h\bigl((g\circ f)(x)\bigr)
=h\bigl(g(f(x))\bigr),\\
\bigl((h\circ g)\circ f\bigr)(x)
&=(h\circ g)\bigl(f(x)\bigr)
=h\bigl(g(f(x))\bigr).
\end{align*}
The two sides are the same value, for every $x$ in the domain. The \emph{identity function} $\operatorname{id}_E$ sends each $x\in E$ to itself.

An \emph{inverse function} for $f:E\to F$ is a function $g:F\to E$ satisfying both
\[
g\circ f=\operatorname{id}_E\qquad\text{and}\qquad f\circ g=\operatorname{id}_F.
\]
The first equation undoes $f$ on inputs. The second undoes $f$ on outputs.

The statement is an ``if and only if'', so both directions have to be written. Uniqueness is a separate claim, with its own step.

\begin{theorem}[the inverse criterion]\label{thm:lang-inverse}
A function has an inverse function if and only if it is bijective. The inverse, when it exists, is unique, and it is denoted $f^{-1}$.
\end{theorem}
\begin{proof}
We prove the two directions, and then uniqueness.

\emph{An inverse forces a bijection.} Suppose $g$ is an inverse of $f$. If $f(x)=f(x')$, apply $g$ to both sides. The identity $g\circ f=\operatorname{id}_E$ gives $x=g(f(x))=g(f(x'))=x'$. So $f$ is injective. For each $y\in F$, the identity $f\circ g=\operatorname{id}_F$ says $f(g(y))=y$. The element $g(y)$ is a preimage of $y$. So $f$ is surjective.

\emph{A bijection produces an inverse.} Suppose $f$ is bijective. Take $y\in F$. Surjectivity supplies at least one $x\in E$ with $f(x)=y$. Injectivity says there is only one such $x$. Define $g(y)$ to be that unique $x$. Then $f(g(y))=y$ by construction, so $f\circ g=\operatorname{id}_F$. For the other identity, take $x\in E$ and set $y=f(x)$. The element $x$ is a preimage of $y$, and it is the only one, so $g(y)=x$. That is $g(f(x))=x$. Thus $g\circ f=\operatorname{id}_E$.

\emph{Uniqueness.} Suppose $h$ is another inverse. For every $y\in F$,
\[
h(y)=h\bigl(f(g(y))\bigr)=g(y),
\]
because $f(g(y))=y$ and $h(f(g(y)))=g(y)$. So $h=g$.
\end{proof}

The criterion is the reason both identities appear in the definition. One of them is not a courtesy copy of the other.

\begin{think}
The maps $f:\N\to\N$, $f(n)=n+1$, and $g:\N\to\N$, with $g(0)=0$ and $g(n)=n-1$ for $n\ge 1$, satisfy $g\circ f=\operatorname{id}_{\N}$. Does $f\circ g=\operatorname{id}_{\N}$? What does the answer say about checking only one of the two inverse identities?
\end{think}

Compute the missing identity at $0$. We have $g(0)=0$, so $f(g(0))=f(0)=1$. The identity function sends $0$ to $0$. Since $1\ne 0$, the composite $f\circ g$ is not $\operatorname{id}_{\N}$. The number $0$ is a value that $f$ never takes, because $n+1\ge 1$ for every $n\in\N$. So $f$ is not surjective, and Theorem~\ref{thm:lang-inverse} says it has no inverse. Checking $g\circ f$ alone would have missed the failure.

For a bijection you can often discover a formula by solving. Take $f:\R\to\R$, $f(x)=3x-2$. The equation $y=3x-2$ rearranges to $x=(y+2)/3$, because subtracting $2$ and multiplying by the inverse of $3$ are legal on $\R$. Define $g(y)=(y+2)/3$. Then
\[
g(f(x))=\frac{(3x-2)+2}{3}=x,\qquad f(g(y))=3\cdot\frac{y+2}{3}-2=y.
\]
Both identities hold, on the stated domains. Solving suggests the formula. The two substitutions are the proof. For a bijection, the inverse image of a set $B$ is also the direct image of $B$ under the inverse function, so the two uses of the symbol $f^{-1}(B)$ agree.

Functions compare inputs by their outputs. The next section compares elements of one set directly, by a relation.

\section{Relations, equivalence classes and partitions}

A \emph{relation} on a set $E$ is a subset $R\subseteq E\times E$. Writing $x\mathrel R y$ means $(x,y)\in R$. Equality, the order $\le$ on $\R$, and ``has the same square as'' on $\R$ are relations, and they behave differently. Three properties have names.
\begin{itemize}
\item \emph{Reflexive}: $x\mathrel R x$ for every $x\in E$.
\item \emph{Symmetric}: $x\mathrel R y$ implies $y\mathrel R x$.
\item \emph{Transitive}: $x\mathrel R y$ and $y\mathrel R z$ imply $x\mathrel R z$.
\end{itemize}

\begin{definitionx}[equivalence relation]
An \emph{equivalence relation} on $E$ is a relation that is reflexive, symmetric and transitive. It is usually written $\sim$.
\end{definitionx}

An equivalence relation records a precise sense in which objects are to be treated as alike. The relation $x\sim y$ defined by $x^2=y^2$ is an equivalence relation on $\R$. Reflexivity uses $x^2=x^2$. Symmetry uses the symmetry of equality. If $x^2=y^2$ and $y^2=z^2$, then $x^2=z^2$, which is transitivity.

By contrast, define $x\sim y$ on $\R$ by $|x-y|<1$. Here $|t|$ denotes $t$ when $t\ge 0$ and $-t$ when $t<0$. The relation is reflexive, because $|x-x|=0<1$. It is symmetric, because $|x-y|=|y-x|$. It is not transitive. We have $|0-3/4|=3/4<1$ and $|3/4-3/2|=3/4<1$, but $|0-3/2|=3/2$, which is not strictly less than $1$. Figure~\ref{fig:lang-transitive} draws the three distances. Nearness of this kind does not survive a chain.

\begin{figure}[ht]
\centering
\fig{ch00-transitive}
\caption{Being strictly less than $1$ apart is not transitive. Each solid step has length $3/4$. The dashed step, from $0$ to $3/2$, has length $3/2$.}
\label{fig:lang-transitive}
\end{figure}

The \emph{equivalence class} of $x$ is $[x]=\{y\in E:y\sim x\}$. For the square relation, $[2]=\{-2,2\}$ and $[0]=\{0\}$. Classes need not have the same size. Figure~\ref{fig:lang-classes} shows three classes on the set $\{-2,-1,0,1,2\}$.

\begin{figure}[ht]
\centering
\fig{ch00-classes}
\caption{Same square, same class. The blocks are nonempty, they cover the five numbers, and distinct blocks do not meet. That is a partition.}
\label{fig:lang-classes}
\end{figure}

A \emph{partition} of $E$ is a collection of nonempty subsets, called its \emph{blocks}, whose union is $E$ and whose distinct blocks are disjoint. The empty collection partitions the empty set. A partition sorts objects into blocks with no omissions and no overlap.

The theorem says these are two descriptions of the same thing. Write each direction on its own.

\begin{theorem}[equivalence relations are partitions]\label{thm:lang-partition}
The distinct equivalence classes of an equivalence relation on $E$ form a partition of $E$. Conversely, any partition defines an equivalence relation by declaring two elements equivalent exactly when they lie in the same block. The classes of that relation are the original blocks.
\end{theorem}
\begin{proof}
Let $\sim$ be an equivalence relation on $E$.

\emph{The classes cover $E$, and they are nonempty.} Reflexivity gives $x\sim x$, so $x\in[x]$. Every class is therefore nonempty, and every element lies in a class.

\emph{Classes that meet are equal.} Suppose $z\in[x]\cap[y]$. Then $z\sim x$ and $z\sim y$. Symmetry turns $z\sim x$ into $x\sim z$. Transitivity with $z\sim y$ gives $x\sim y$. Now take any $t\in[x]$, so $t\sim x$. Transitivity with $x\sim y$ gives $t\sim y$, hence $t\in[y]$. This shows $[x]\subseteq[y]$. The same argument with $x$ and $y$ exchanged shows $[y]\subseteq[x]$. So $[x]=[y]$. Distinct classes are therefore disjoint: if they met, they would be equal.

\emph{From a partition back to a relation.} Let a partition of $E$ be given, and declare $x\sim y$ when $x$ and $y$ lie in the same block. Every element lies in exactly one block, so $x\sim x$. Sharing a block does not depend on order, so the relation is symmetric. Suppose $x\sim y$ and $y\sim z$. Then $x$ and $y$ share a block, and $y$ and $z$ share a block. Those two blocks both contain $y$, so they are the same block, because distinct blocks are disjoint. Thus $x$ and $z$ share a block, and $x\sim z$. The elements related to $x$ are exactly the elements of its block. So the classes are the original blocks.
\end{proof}

On the five-element set of Figure~\ref{fig:lang-classes}, the three blocks are the classes of the square relation, restricted to that set. They cover every listed number, and no number sits in two blocks. The theorem says this is the general picture, not a lucky drawing.

\begin{beauty}
One relation sorts the set with no overlaps and no leftovers, and the sorting gives the relation back. Remainder classes, the fibres of a function, and the outcomes a measurement cannot tell apart are this same theorem again, later. The definition already contains the argument.
\end{beauty}

The same idea of comparison can refuse to treat distinct objects as alike. That is an order.

\section{Order: comparable does not mean equal}

A \emph{partial order} on a set $E$ is a reflexive, transitive relation $\preceq$ which is also \emph{antisymmetric}: if $x\preceq y$ and $y\preceq x$, then $x=y$. Antisymmetry and symmetry ask for different things. In an order, comparison both ways forces equality. In an equivalence relation, comparison both ways is expected even when the objects are distinct. The usual relation $\le$ on $\R$ is a partial order. So is inclusion on $\mathcal P(E)$: every set contains itself, inclusions compose, and two opposite inclusions give equality of sets.

An order is \emph{total} when any two elements are comparable, meaning $x\preceq y$ or $y\preceq x$. Inclusion is usually only partial. The sets $\{1\}$ and $\{2\}$ are incomparable: neither contains the other. The associated strict order $x\prec y$ means $x\preceq y$ and $x\ne y$.

Let $A$ be a subset of an ordered set. A \emph{least element} of $A$ is an element $a\in A$ such that $a\preceq x$ for every $x\in A$. A \emph{minimal element} of $A$ is an element of $A$ with no strictly smaller element of $A$. We compare the two words in steps.

\emph{Least implies minimal.} Suppose $a$ is least in $A$, and suppose $x\in A$ satisfies $x\prec a$. Then $x\preceq a$ and $x\ne a$. Leastness also gives $a\preceq x$. Antisymmetry forces $x=a$, which contradicts $x\ne a$. So no such $x$ exists, and $a$ is minimal.

\emph{The converse can fail.} Consider $\mathcal A=\{\{1\},\{2\},\{1,2\}\}$, ordered by inclusion. Figure~\ref{fig:lang-order} draws it. The sets $\{1\}$ and $\{2\}$ are minimal: nothing in $\mathcal A$ is strictly contained in either of them. Neither is least, because neither is contained in the other. The set $\{1,2\}$ is a greatest element of $\mathcal A$, because it contains every member.

\emph{In a total order, minimal implies least.} Suppose $a$ is minimal in $A$, and the order is total. Take any $x\in A$. Totality gives $a\preceq x$ or $x\preceq a$. If $x\preceq a$ and $x\ne a$, then $x\prec a$, which contradicts minimality. So that alternative forces $x=a$, and then $a\preceq x$. The other alternative is already $a\preceq x$. In both cases $a\preceq x$. Since $x$ was arbitrary, $a$ is least.

A least element, when it exists, is unique. If $a$ and $b$ are both least, then $a\preceq b$ and $b\preceq a$, so antisymmetry gives $a=b$.

\begin{figure}[ht]
\centering
\fig{ch00-order}
\caption{Under inclusion, $\{1\}$ and $\{2\}$ are minimal and incomparable, so neither is least. The set $\{1,2\}$ is greatest. The empty set, drawn dashed, is a lower bound in the larger power set. It is not an element of this three-element ordered set.}
\label{fig:lang-order}
\end{figure}

Replacing ``below'' by ``above'' defines a \emph{greatest} element and a \emph{maximal} element. An \emph{upper bound} for $A$ is an element $u$ of the ambient ordered set such that $x\preceq u$ for every $x\in A$. The bound need not belong to $A$. A \emph{least upper bound}, or \emph{supremum}, is the least of the upper bounds. Reversing the order defines a \emph{lower bound} and a \emph{greatest lower bound}, or \emph{infimum}.

The open interval $(0,1)\subseteq\R$ has supremum $1$, and it has no greatest element. The proof is a useful rehearsal of cases.

The number $1$ is an upper bound: if $0<x<1$, then $x\le 1$. Now suppose $u<1$. We produce a point of $(0,1)$ that is strictly larger than $u$, so $u$ is not an upper bound. If $u\le 0$, the point $1/2$ lies in $(0,1)$ and satisfies $1/2>u$. If $0<u<1$, set $m=(u+1)/2$. Then
\[
m-u=\frac{1-u}{2}>0,\qquad 1-m=\frac{1-u}{2}>0,
\]
so $u<m<1$. Also $m>0$. Thus $m\in(0,1)$ and $m>u$. Every number strictly below $1$ fails to be an upper bound. Combined with the fact that $1$ is an upper bound, this says that $1$ is the least upper bound.

The same point $m$ shows there is no greatest element. If $g\in(0,1)$, the number $(g+1)/2$ lies in $(0,1)$ and is strictly larger than $g$. So $g$ is not an upper bound of the interval.

Completeness, assumed in the preface, says that every nonempty set of reals that is bounded above has a supremum in $\R$. The interval $(0,1)$ is one example, and its supremum was found by hand. Chapter~\ref{ch:sequences} shows how the same assumption produces limits and existence theorems. The order language above is what that chapter will use.

The workshop uses these words on one question: when does an image preserve an intersection?

\section{Proof workshop: when does an image preserve an intersection?}

The question
\[
f(A\cap B)\stackrel{?}{=}f(A)\cap f(B)
\]
joins definitions, quantifiers, a counterexample and a repair. Here $f:E\to F$ is a function and $A,B\subseteq E$. An element of the left side comes from one input that belongs to both sets. An element of the right side comes from an input in $A$ and an input in $B$. Those two inputs need not be the same element. That is where to experiment.

\subsection*{A worked experiment and a conjecture}

Take $f:\R\to\R$, $f(x)=x^2$, and choose $A=\{-1\}$ and $B=\{1\}$. Then $A\cap B=\varnothing$, so $f(A\cap B)=\varnothing$. But $f(A)=\{1\}$ and $f(B)=\{1\}$, so $f(A)\cap f(B)=\{1\}$. Equality fails. Figure~\ref{fig:lang-intersection} draws the same failure with the names $x$ and $x'$: two inputs, one output, and an empty intersection.

\begin{figure}[ht]
\centering
\fig{ch00-intersection}
\caption{The image of an intersection can be smaller than the intersection of the images. The inputs $x$ and $x'$ are different, so $A\cap B$ is empty, but both are sent to $y$.}
\label{fig:lang-intersection}
\end{figure}

One inclusion seems to survive. Any input that lies in both sets produces an output that lies in both images. The failure used two different inputs with the same output. This suggests that injectivity may be the missing hypothesis, and that it may be exactly the right one.

\begin{think}
Prove the inclusion that always holds. Then suppose $f$ is injective, and try to reverse the inclusion. Which two witnesses does membership in $f(A)\cap f(B)$ supply, and what lets you identify them?
\end{think}

\subsection*{Mission and staged hints}

Prove both assertions.
\begin{enumerate}[label=(\alph*)]
\item For every function $f:E\to F$ and every $A,B\subseteq E$, one has $f(A\cap B)\subseteq f(A)\cap f(B)$. Equality holds if $f$ is injective.
\item If equality holds for every pair of subsets $A,B\subseteq E$, then $f$ is injective.
\end{enumerate}
\emph{First hint.} Start with an arbitrary element of the set whose inclusion you want, and translate ``image'' into a witness.
\emph{Second hint.} For (b), you may choose the sets. If $f(x)=f(x')$, what do $A=\{x\}$ and $B=\{x'\}$ reveal?

\subsection*{How to find it}

The first direction follows the definitions forwards. The failed reverse direction reveals a specific obstruction: the two witnesses can disagree. Injectivity removes that obstruction. To show that the hypothesis is necessary, make the obstruction as small as possible by using singleton sets. This is a common way to find a proof. Attempt the argument, locate the exact failing step, and then either repair the hypothesis or turn the failure into a counterexample.

\subsection*{Write-up}

Let $y\in f(A\cap B)$. By the definition of image, there is an element $x\in A\cap B$ with $f(x)=y$. This $x$ belongs to $A$ and to $B$, so $y\in f(A)$ and $y\in f(B)$. Thus $y\in f(A)\cap f(B)$. This proves the inclusion that needs no extra hypothesis.

Now assume $f$ is injective, and let $y\in f(A)\cap f(B)$. There is an element $a\in A$ with $f(a)=y$, and there is an element $b\in B$ with $f(b)=y$. Then $f(a)=f(b)$. Injectivity gives $a=b$. That common element lies in $A\cap B$, and its image is $y$, so $y\in f(A\cap B)$. Together with the first inclusion, this is equality.

Conversely, assume equality for every pair of subsets of $E$. Suppose $f(x)=f(x')$. Call this common value $y$, and set $A=\{x\}$ and $B=\{x'\}$. Then $y\in f(A)$ and $y\in f(B)$, so $y\in f(A)\cap f(B)$. The assumed equality puts $y$ in $f(A\cap B)$, so $A\cap B$ is nonempty. Two singleton sets meet only when their elements are equal. Thus $x=x'$. This is injectivity.

If $E=\varnothing$, there are no distinct inputs to separate, and every inclusion of subsets of the empty set is an equality of empty sets. The injectivity condition and the inclusions hold vacuously.

\subsection*{A tempting invalid proof}

``If $y\in f(A)\cap f(B)$, write $y=f(x)$ with $x\in A$ and $x\in B$.'' The single letter $x$ has identified two independent existential choices. The correct translation begins with $a\in A$ and $b\in B$. Give different witnesses different names until an argument proves that they coincide. One letter for two choices is the same slip as misreading the order of the quantifiers: a choice made twice is two choices, until a reason makes them one.

\section*{What this chapter prepares}
You can now hear what a sentence is asking. At the start, $x^2=1$ asked three different things, depending on the domain and the words around it. A proof starts from those words: a witness, a choice that may depend on an earlier choice, two inclusions, both directions.

The next chapter asks whether
\[
1^3+2^3+\dots+n^3=(1+2+\dots+n)^2
\]
holds for every positive integer $n$, or only for the values already checked.

\section{Exercises}

\begin{exercise}[$\star$]\label{ex:lang-implications}
For a real number $x$, decide which implications hold for every $x$, and give a counterexample to each that fails:
(a) $x=2\Rightarrow x^2=4$;
(b) $x^2=4\Rightarrow x=2$;
(c) $x>2\Rightarrow x^2>4$;
(d) $x^2>4\Rightarrow x>2$.
Give a necessary and sufficient condition for $x^2>4$.
\end{exercise}

\begin{exercise}[$\star$]\label{ex:lang-vacuous}
Decide whether each sentence is true.
(a) Every integer that is both even and odd is greater than $100$.
(b) Some integer that is both even and odd is greater than $100$.
(c) For every real $x$, one has $x>0$ or $x<1$. Name a real number at which both parts are true.
\end{exercise}

\begin{exercise}[$\star$]\label{ex:lang-negation}
Negate each statement, and keep its domain:
(a) every $x\in\R$ satisfies $x^2\ge x$;
(b) some $x\in\R$ satisfies $0<x<1$ and $x^2=x$;
(c) for every $x\in\R$ there is $y\in\R$ with $y^2=x$.
Decide whether the original statements are true.
\end{exercise}

\begin{exercise}[$\star$]\label{ex:lang-quantifiers}
Decide which statements are true, and justify the answers:
\[
\forall x\in\R\ \exists y\in\R:\ y>x,
\qquad
\exists y\in\R\ \forall x\in\R:\ y>x.
\]
Write their negations, and explain which choices may depend on which variables.
\end{exercise}

\begin{exercise}[$\star$]\label{ex:lang-setalgebra}
Prove $A\setminus(B\cup C)=(A\setminus B)\cap(A\setminus C)$ by comparing membership. Is $A\setminus(B\cap C)=(A\setminus B)\cap(A\setminus C)$ always true? If it fails, give a counterexample and a corrected equality.
\end{exercise}

\begin{exercise}[$\star$]\label{ex:lang-powerset}
List $\mathcal P(\{\varnothing,\{1\}\})$. Which of the following hold for an arbitrary set $A$:
$\varnothing\in\mathcal P(A)$, $\varnothing\subseteq\mathcal P(A)$, $A\in\mathcal P(A)$, and $A\subseteq\mathcal P(A)$?
Give a counterexample to any assertion which need not hold.
\end{exercise}

\begin{exercise}[$\star\star$]\label{ex:lang-domain}
Let $E=\R\setminus\{0\}$, and let $f:E\to\R$ be given by $f(x)=1/x$. Prove that $f$ is injective and not surjective. Change only the codomain so that the same rule becomes bijective, and find and verify the inverse. Explain why the codomain matters.
\end{exercise}

\begin{exercise}[$\star\star$]\label{ex:lang-inverseimage}
For a function $f:E\to F$ and sets $A\subseteq E$, $B\subseteq F$, prove
\[
A\subseteq f^{-1}(f(A)),\qquad f(f^{-1}(B))=B\cap f(E).
\]
Give an example where the first inclusion is strict. Characterise the functions for which that inclusion is an equality for every $A\subseteq E$.
\end{exercise}

\begin{exercise}[$\star\star$]\label{ex:lang-complement}
For $f:E\to F$, prove $f(E\setminus A)\supseteq f(E)\setminus f(A)$. Give a counterexample to equality. Prove equality when $f$ is injective. Explain why replacing $f(E)$ by $F$ would also require surjectivity.
\end{exercise}

\begin{exercise}[$\star\star$]\label{ex:lang-composition}
Let $f:E\to F$ and $g:F\to G$. Prove:
(a) if both are injective, then $g\circ f$ is injective;
(b) if both are surjective, then $g\circ f$ is surjective;
(c) if $g\circ f$ is injective, then $f$ is injective;
(d) if $g\circ f$ is surjective, then $g$ is surjective.
Give examples showing that the other factor need not have the same property in (c) or in (d).
\end{exercise}

\begin{exercise}[$\star\star$]\label{ex:lang-selfinverse}
Let $E=\R\setminus\{1\}$. Show that $f(x)=(x+1)/(x-1)$ defines a function $E\to E$ and satisfies $f\circ f=\operatorname{id}_E$. Deduce bijectivity. Check every denominator before you cancel it.
\end{exercise}

\begin{exercise}[$\star\star$]\label{ex:lang-equivalence}
Define $x\sim y$ on $\R$ by $x-y\in\Z$. Prove that this is an equivalence relation. Describe $[0]$ and $[1/2]$, and decide whether $[1/2]=[3/2]$. Can two distinct equivalence classes share an element?
\end{exercise}

\begin{exercise}[$\star\star$]\label{ex:lang-partitions}
Suppose $f:E\to F$. Define $x\sim y$ when $f(x)=f(y)$. Prove that $\sim$ is an equivalence relation, and that its classes are exactly the nonempty fibres $f^{-1}(\{y\})$ for $y\in f(E)$. What does injectivity say about these classes?
\end{exercise}

\begin{exercise}[$\star\star$]\label{ex:lang-order}
Order $\mathcal A=\{\{1\},\{2\},\{1,2\}\}$ by inclusion, as in Figure~\ref{fig:lang-order}. Prove that $\{1\}$ and $\{2\}$ are minimal, that $\mathcal A$ has no least element, and that $\{1,2\}$ is the greatest element. Prove that any least element of any partially ordered set is unique. In the larger ordered set $\mathcal P(\{1,2\})$, is $\varnothing$ an upper bound for $\mathcal A$, or a lower bound?
\end{exercise}

\begin{exercise}[$\star\star$]\label{ex:lang-rational}
Let $r\in\Q$ and $x\in\R$. Prove that $x+r$ is rational if and only if $x$ is rational. If also $r\ne 0$, prove that $rx$ is rational if and only if $x$ is rational. State the corresponding consequences for irrational numbers, and explain the role of $r\ne 0$.
\end{exercise}

\begin{exercise}[$\star\star$]\label{ex:lang-existence}
For each $a\in(0,1)$, prove that there is a unique $x>0$ satisfying $x/(1+x)=a$. Explain separately how the assumptions give an admissible candidate, and how the equation gives uniqueness.
\end{exercise}

\begin{exercise}[$\star$]\label{ex:lang-falseproof}
Find the first invalid step in this purported proof that $1=2$. Let $a=b=1$. Then $a^2=ab$, hence $a^2-b^2=ab-b^2$, hence $(a-b)(a+b)=b(a-b)$, hence $a+b=b$, hence $2=1$. Why do the valid lines before the error fail to make the conclusion plausible?
\end{exercise}

\begin{exercise}[$\star\star$]\label{ex:lang-midpoint}
For real numbers $a<b$, prove that there is exactly one real $x$ with $|x-a|=|x-b|$, and determine it. Give one proof by cases on the position of $x$, and another by squaring. Explain why squaring is reversible in this argument.
\end{exercise}

\begin{exercise}[$\star\star\star$]\label{ex:lang-retraction}
Let $f:E\to E$ satisfy $f\circ f=f$.
(a) Prove $f(E)=\{y\in E:f(y)=y\}$.
(b) Prove that each class of the relation $x\sim x'$ defined by $f(x)=f(x')$ contains exactly one fixed point of $f$.
(c) If such an $f$ is injective, or if it is surjective, prove $f=\operatorname{id}_E$.
Give an example with $f\ne\operatorname{id}_E$ when neither property is required.
\end{exercise}

\section{Hints}

\begin{hint}{ex:lang-implications}
Test the negative root in (b), and a number below $-2$ in (d). Factor $x^2-4=(x-2)(x+2)$.
\end{hint}
\begin{hint}{ex:lang-vacuous}
For (a) and (b), ask whether any integer is both even and odd. For (c), split into the cases $x\le 0$ and $x>0$, and remember that mathematical ``or'' allows both parts.
\end{hint}
\begin{hint}{ex:lang-negation}
Change each quantifier, then negate its predicate. The negation of an ``and'' is an ``or''.
\end{hint}
\begin{hint}{ex:lang-quantifiers}
For a given $x$, try $y=x+1$. For a proposed $y$ that is supposed to work for every $x$, choose an $x$ from that $y$.
\end{hint}
\begin{hint}{ex:lang-setalgebra}
Use the phrases ``in $A$, and outside both $B$ and $C$'' and ``in $A$, and outside at least one of $B$ and $C$''. Test sets with one element.
\end{hint}
\begin{hint}{ex:lang-powerset}
The base set has two elements, even though one of them is empty. Membership in $\mathcal P(A)$ means being a subset of $A$.
\end{hint}
\begin{hint}{ex:lang-domain}
Which real value can a reciprocal never take? To solve $1/x=y$, you need $y\ne 0$.
\end{hint}
\begin{hint}{ex:lang-inverseimage}
Translate each membership into a witness. For the characterisation, use a singleton set $A$.
\end{hint}
\begin{hint}{ex:lang-complement}
A value in $f(E)$ but not in $f(A)$ has a preimage, and none of its preimages lies in $A$. For a strict inclusion, let two different inputs share a value.
\end{hint}
\begin{hint}{ex:lang-composition}
Write out each definition. For a failed converse, a small intermediate set can hide bad behaviour elsewhere in $F$.
\end{hint}
\begin{hint}{ex:lang-selfinverse}
First show that $(x+1)/(x-1)=1$ is impossible. Then compute $f(x)-1$ and $f(x)+1$ over the common denominator $x-1$.
\end{hint}
\begin{hint}{ex:lang-equivalence}
Use $0\in\Z$, and use closure of $\Z$ under negation and under addition. Each class is a set of integer translates.
\end{hint}
\begin{hint}{ex:lang-partitions}
Equality of outputs inherits reflexivity, symmetry and transitivity from equality of real numbers, or of elements of $F$. Fix the output of one input.
\end{hint}
\begin{hint}{ex:lang-order}
Draw the three sets, with a line upward for strict inclusion. A bound in the larger ordered set need not belong to $\mathcal A$. Figure~\ref{fig:lang-order} is the drawing.
\end{hint}
\begin{hint}{ex:lang-rational}
Use closure of the rational numbers under addition and subtraction, and under division by a nonzero rational. Write the fraction arithmetic if you need to see the closure.
\end{hint}
\begin{hint}{ex:lang-existence}
Solve the equation backwards, then check the candidate forwards. The sign of $1-a$ decides whether the candidate is positive.
\end{hint}
\begin{hint}{ex:lang-falseproof}
Which factor was cancelled, and what is its value when $a=b=1$?
\end{hint}
\begin{hint}{ex:lang-midpoint}
Outside $[a,b]$, one distance is strictly larger than the other. Inside, the equation is linear. For the second proof, expand the difference of the two squares.
\end{hint}
\begin{hint}{ex:lang-retraction}
Apply $f\circ f=f$ to an element of the image. In the class of $x$, try the representative $f(x)$.
\end{hint}

\section{Solutions}

\begin{solution}{ex:lang-implications}
\idea Squares forget the sign of a number. A permitted input is not the same thing as an arbitrary input.
\writeup (a) holds for every real $x$. If $x=2$, then $x^2=4$ by substitution.

(b) fails at $x=-2$. There $(-2)^2=4$, and $-2\ne 2$.

(c) holds for every real $x$. Suppose $x>2$. Then $x-2>0$ and $x+2>2+2=4>0$, so
\[
x^2-4=(x-2)(x+2)>0.
\]
Thus $x^2>4$.

(d) fails at $x=-3$. There $(-3)^2=9>4$, and $-3>2$ is false.

The necessary and sufficient condition is $x<-2$ or $x>2$. First suppose $x<-2$ or $x>2$. In the second sub-case, (c) gives $x^2>4$. In the first, $x+2<0$ and $x-2<x+2<0$, so both factors of $x^2-4$ are negative, and their product is positive. Thus $x^2>4$. Conversely, suppose $-2\le x\le 2$. Then $|x|\le 2$. The square function is increasing on $[0,\infty)$, in the sense that $0\le u\le v$ implies $u^2\le v^2$, because $v^2-u^2=(v-u)(v+u)$ and both factors are nonnegative. Apply this to $u=|x|$ and $v=2$ to get $x^2\le 4$. So $x^2>4$ fails throughout $[-2,2]$.
\end{solution}

\begin{solution}{ex:lang-vacuous}
\idea An empty collection makes a universal sentence true and an existential sentence false. Inclusive ``or'' stays true when both parts hold.
\writeup (a) No integer is both even and odd. Suppose $n$ is even and $n$ is odd. By definition, $n=2k$ and $n=2m+1$ for integers $k$ and $m$. Then $2k-2m=1$, so $2(k-m)=1$. The quantifier section already showed that twice an integer never equals $1$: if the integer factor is at least $1$, the product is at least $2$, and if it is at most $0$, the product is at most $0$. So the hypothesis ``$n$ is both even and odd'' never holds, and it supplies no counterexample to ``$n>100$''. The universal sentence is true. This is vacuous truth.

(b) The existential sentence needs a witness in $\Z$. The argument in (a) shows there is none. So (b) is false.

(c) Take a real $x$. If $x\le 0$, then $x<1$, so the second part holds. If $x>0$, the first part holds. In both cases the inclusive or holds, so the universal sentence is true. At $x=1/2$ both parts are true: $1/2>0$ and $1/2<1$.
\end{solution}

\begin{solution}{ex:lang-negation}
\idea Keep the domain fixed. Negate from the outside inward, using Proposition~\ref{prop:lang-negation}. Then test a simple candidate.
\writeup (a) The original statement is $\forall x\in\R,\ x^2\ge x$. Its negation is $\exists x\in\R,\ x^2<x$. The original is false, because at $x=1/2$ one has $(1/2)^2=1/4<1/2$.

(b) The original statement is $\exists x\in\R,\ (0<x<1\ \text{and}\ x^2=x)$. Its negation is
\[
\forall x\in\R,\quad x\le 0\ \text{or}\ x\ge 1\ \text{or}\ x^2\ne x.
\]
The original is false. The equation $x^2=x$ rearranges to $x^2-x=0$, hence $x(x-1)=0$, so $x=0$ or $x=1$. Both of these fail $0<x<1$. There is no witness in the open interval.

(c) The original statement is $\forall x\in\R\ \exists y\in\R,\ y^2=x$. Its negation is $\exists x\in\R\ \forall y\in\R,\ y^2\ne x$. The original is false, because $x=-1$ is not a square of a real number: for every real $y$, one has $y^2\ge 0$, so $y^2\ne -1$.
\end{solution}

\begin{solution}{ex:lang-quantifiers}
\idea Ask what information is available at the moment $y$ is chosen.
\writeup The first statement is true. Given $x\in\R$, choose $y=x+1$. Then $y-x=1>0$, so $y>x$. Here $y$ may depend on $x$.

The negation of the first statement, by Proposition~\ref{prop:lang-negation} applied twice from the outside, is
\[
\exists x\in\R\ \forall y\in\R,\quad y\le x.
\]
That sentence is false, which agrees with the first statement being true.

The second statement is false. Suppose some $y\in\R$ satisfied $y>x$ for every real $x$. Taking $x=y$ would require $y>y$, which is false. So no such $y$ exists. In this sentence the witness $y$ is chosen before $x$, and it is not allowed to depend on $x$.

The negation of the second statement is $\forall y\in\R\ \exists x\in\R,\ y\le x$. This is true: given $y$, choose $x=y$. Then $y\le x$.
\end{solution}

\begin{solution}{ex:lang-setalgebra}
\idea Difference means membership together with non-membership. The two sides of a proposed equality can be tested on a one-element set.
\writeup Let $x$ be arbitrary. Then $x\in A\setminus(B\cup C)$ means $x\in A$, and $x\notin B\cup C$. The second demand means that $x\notin B$ and $x\notin C$. So $x\in A\setminus B$ and $x\in A\setminus C$, which means $x\in(A\setminus B)\cap(A\setminus C)$. Each step reverses, so the two sets are equal.

The proposed equality $A\setminus(B\cap C)=(A\setminus B)\cap(A\setminus C)$ is not always true. Take $A=B=\{1\}$ and $C=\varnothing$. The left side is $\{1\}\setminus(\{1\}\cap\varnothing)=\{1\}\setminus\varnothing=\{1\}$. The right side is $(\{1\}\setminus\{1\})\cap(\{1\}\setminus\varnothing)=\varnothing\cap\{1\}=\varnothing$. The two sides differ.

The corrected equality is $A\setminus(B\cap C)=(A\setminus B)\cup(A\setminus C)$. Indeed, if $x\in A$, then $x\notin B\cap C$ means that $x\notin B$ or $x\notin C$, by Proposition~\ref{prop:lang-negation}. That is $x\in A\setminus B$ or $x\in A\setminus C$.
\end{solution}

\begin{solution}{ex:lang-powerset}
\idea Name the two elements of the base set before you ask what they contain.
\writeup The set $A_0=\{\varnothing,\{1\}\}$ has two distinct elements, $\varnothing$ and $\{1\}$. Its subsets are obtained by including each element or not:
\[
\mathcal P(A_0)=\bigl\{\varnothing,\ \{\varnothing\},\ \{\{1\}\},\ \{\varnothing,\{1\}\}\bigr\}.
\]

For an arbitrary set $A$, the first three assertions hold. The empty set is a subset of $A$, so $\varnothing\in\mathcal P(A)$. The empty set is a subset of every set, so $\varnothing\subseteq\mathcal P(A)$. And $A\subseteq A$, so $A\in\mathcal P(A)$.

The inclusion $A\subseteq\mathcal P(A)$ need not hold. It asks every element of $A$ to be a subset of $A$. Take $A=\{\{\varnothing\}\}$. Its only element is $\{\varnothing\}$. That element is a subset of $A$ only if $\varnothing\in A$. But the only element of $A$ is $\{\varnothing\}$, and $\varnothing\ne\{\varnothing\}$, so $\varnothing\notin A$. Thus $\{\varnothing\}\not\subseteq A$, and $A\not\subseteq\mathcal P(A)$.
\end{solution}

\begin{solution}{ex:lang-domain}
\idea Equal outputs give an equation you can clear. The missing outputs are the values the formula never hits.
\writeup Suppose $1/x=1/x'$ with $x,x'\in E$, so $x\ne 0$ and $x'\ne 0$. Multiply both sides by $xx'$, which is legal and nonzero. The left side becomes $x'$, and the right side becomes $x$. Thus $x=x'$. The function is injective.

The value $0$ is never attained: $1/x=0$ would give $1=0$ after multiplying by $x$. So $f:E\to\R$ is not surjective.

Keep the same rule and change the codomain to $E$. For every $y\in E$ one has $y\ne 0$, so $x=1/y$ lies in $E$, and $f(x)=1/(1/y)=y$. The same calculation shows that the inverse function is $f$ itself: $f(f(x))=x$ for every $x\in E$. Both inverse identities hold because the two maps are the same map. By Theorem~\ref{thm:lang-inverse}, the function $E\to E$ is bijective.

The codomain decides surjectivity. The rule did not change. The set of required outputs did.
\end{solution}

\begin{solution}{ex:lang-inverseimage}
\idea A forward image followed by a preimage can collect extra inputs that share an output.
\writeup Suppose $x\in A$. Then $f(x)\in f(A)$ by definition of image, so $x\in f^{-1}(f(A))$ by definition of preimage. This is the first inclusion.

Suppose $y\in f(f^{-1}(B))$. Then $y=f(x)$ for some $x\in f^{-1}(B)$. For that $x$ one has $f(x)\in B$, so $y\in B$. Also $y=f(x)$ lies in $f(E)$. Thus $y\in B\cap f(E)$. Conversely, suppose $y\in B\cap f(E)$. Then $y=f(x)$ for some $x\in E$, and $f(x)\in B$, so $x\in f^{-1}(B)$. Therefore $y\in f(f^{-1}(B))$. The two inclusions give the stated equality.

For a strict inclusion, take $f:\R\to\R$, $f(x)=x^2$, and $A=\{1\}$. Then $f(A)=\{1\}$ and $f^{-1}(f(A))=\{-1,1\}$, which properly contains $A$.

If $f$ is injective and $x\in f^{-1}(f(A))$, choose $a\in A$ with $f(x)=f(a)$. Injectivity gives $x=a$, so $x\in A$. Combined with the first inclusion, one has equality for every $A$.

Conversely, suppose equality holds for every $A\subseteq E$, and suppose $f(x)=f(x')$. Set $A=\{x\}$. Then $f(x')=f(x)\in f(A)$, so $x'\in f^{-1}(f(A))=A=\{x\}$. Thus $x'=x$. The exact condition is injectivity.
\end{solution}

\begin{solution}{ex:lang-complement}
\idea An output missing from $f(A)$ cannot have a preimage in $A$. An output that does lie in $f(A)$ may still have another preimage outside $A$.
\writeup Suppose $y\in f(E)\setminus f(A)$. Choose $x\in E$ with $f(x)=y$. If $x$ were in $A$, then $y$ would lie in $f(A)$, contrary to the choice of $y$. So $x\in E\setminus A$, and therefore $y\in f(E\setminus A)$. This is the inclusion.

Equality can fail. Take $f(x)=x^2$ on $\R$ and $A=\{1\}$. Then $f(-1)=1$ and $-1\notin A$, so $1\in f(E\setminus A)$. But $1=f(1)$ lies in $f(A)$, so $1\notin f(E)\setminus f(A)$.

Now assume $f$ is injective, and take $y\in f(E\setminus A)$. Write $y=f(x)$ with $x\notin A$. If there were an $a\in A$ with $f(a)=y$, injectivity would give $a=x$, contradicting $x\notin A$. So $y\notin f(A)$. Also $y\in f(E)$. Thus $y\in f(E)\setminus f(A)$. This is the reverse inclusion.

Finally, suppose one wants $f(E\setminus A)=F\setminus f(A)$ for every $A\subseteq E$. Take $A=\varnothing$. Then $f(E\setminus\varnothing)=f(E)$ and $F\setminus f(\varnothing)=F\setminus\varnothing=F$, so the equality says $f(E)=F$. That is surjectivity. If $f$ is both injective and surjective, the equality $f(E\setminus A)=F\setminus f(A)$ follows from the equality already proved, because $f(E)=F$.
\end{solution}

\begin{solution}{ex:lang-composition}
\idea Trace equal outputs backwards for injectivity. Trace a desired output backwards for surjectivity.
\writeup (a) Suppose $(g\circ f)(x)=(g\circ f)(x')$. Then $g(f(x))=g(f(x'))$. Injectivity of $g$ gives $f(x)=f(x')$. Injectivity of $f$ gives $x=x'$.

(b) Take $z\in G$. Surjectivity of $g$ supplies $y\in F$ with $g(y)=z$. Surjectivity of $f$ supplies $x\in E$ with $f(x)=y$. Then $(g\circ f)(x)=z$.

(c) Suppose $f(x)=f(x')$. Apply $g$ to get $g(f(x))=g(f(x'))$. Injectivity of $g\circ f$ gives $x=x'$.

(d) Take $z\in G$. Surjectivity of $g\circ f$ supplies $x\in E$ with $g(f(x))=z$. The element $f(x)$ is a preimage of $z$ under $g$.

For both failed converses, take $E=G=\{0\}$, $F=\{0,1\}$, $f(0)=0$ and $g(0)=g(1)=0$. Then $g\circ f$ sends $0$ to $0$, so it is a bijection of a one-element set. In particular it is injective and surjective. But $g$ is not injective, because $g(0)=g(1)$ and $0\ne 1$. And $f$ is not surjective onto $F$, because $1$ is missed. So in (c) the outer map need not be injective, and in (d) the inner map need not be surjective.
\end{solution}

\begin{solution}{ex:lang-selfinverse}
\idea Show that the second application is legal before you simplify it.
\writeup Take $x\in E$, so $x\ne 1$. The denominator $x-1$ is nonzero, and $f(x)=(x+1)/(x-1)$ is a real number.

This value is not $1$. If $(x+1)/(x-1)=1$, then $x+1=x-1$, because the denominator is nonzero and we may multiply by it. Then $1=-1$, which is false. So $f(x)\ne 1$, and $f(x)\in E$. The rule defines a function $E\to E$.

Now compute, still with $x\ne 1$,
\[
f(x)+1=\frac{x+1}{x-1}+\frac{x-1}{x-1}=\frac{2x}{x-1},
\qquad
f(x)-1=\frac{x+1}{x-1}-\frac{x-1}{x-1}=\frac{2}{x-1}.
\]
The second numerator is $2\ne 0$, so $f(x)-1\ne 0$. Therefore $f(f(x))$ is defined, and
\[
f(f(x))=\frac{f(x)+1}{f(x)-1}
=\frac{2x/(x-1)}{2/(x-1)}
=\frac{2x}{x-1}\cdot\frac{x-1}{2}
=x.
\]
The cancellation of $x-1$ is legal because $x\ne 1$. Thus $f\circ f=\operatorname{id}_E$. Both inverse identities hold, because the two maps in each composite are this same function $f$. Theorem~\ref{thm:lang-inverse} gives bijectivity.
\end{solution}

\begin{solution}{ex:lang-equivalence}
\idea Differences add along a chain, and the integers are closed under that addition.
\writeup For every real $x$, one has $x-x=0\in\Z$, so the relation is reflexive. If $x-y\in\Z$, then $y-x=-(x-y)\in\Z$, because the negative of an integer is an integer. So the relation is symmetric. If $x-y\in\Z$ and $y-z\in\Z$, then
\[
x-z=(x-y)+(y-z)\in\Z,
\]
because the sum of two integers is an integer. So the relation is transitive. It is an equivalence relation.

The class of $0$ is $[0]=\{y\in\R:y-0\in\Z\}=\Z$. The class of $1/2$ is
\[
[1/2]=\{1/2+k:k\in\Z\}.
\]
Since $3/2-1/2=1\in\Z$, one has $3/2\sim 1/2$, so $[1/2]=[3/2]$.

Two classes that share an element are equal, by Theorem~\ref{thm:lang-partition}. Distinct classes therefore do not meet.
\end{solution}

\begin{solution}{ex:lang-partitions}
\idea A fibre records exactly the inputs that a function cannot tell apart.
\writeup Reflexivity: $f(x)=f(x)$. Symmetry: if $f(x)=f(y)$, then $f(y)=f(x)$. Transitivity: if $f(x)=f(y)$ and $f(y)=f(z)$, then $f(x)=f(z)$. So $\sim$ is an equivalence relation.

Take $x\in E$. Its class is $\{x'\in E:f(x')=f(x)\}$, which is the fibre $f^{-1}(\{f(x)\})$. This fibre is nonempty because it contains $x$, and $f(x)\in f(E)$. Conversely, if $y\in f(E)$, then $y=f(x)$ for some $x$, and the fibre $f^{-1}(\{y\})$ is the class of that $x$.

Injectivity holds exactly when every one of these classes has a single element. Each class contains its defining input. A second element would be a second input with the same output. If $E=\varnothing$, there are no classes, and the injectivity condition holds vacuously.
\end{solution}

\begin{solution}{ex:lang-order}
\idea A least element has to lie below every other element. Minimality only forbids a strictly smaller one. Figure~\ref{fig:lang-order} is the picture.
\writeup Under inclusion, $\{1\}$ and $\{2\}$ are minimal: no member of $\mathcal A$ is strictly contained in either of them. Neither is contained in the other, so there is no least element. The set $\{1,2\}$ contains every member of $\mathcal A$, so it is the greatest element, and in particular it is the only maximal element.

If $a$ and $b$ are both least in any partial order, then $a\preceq b$ and $b\preceq a$. Antisymmetry gives $a=b$.

In $\mathcal P(\{1,2\})$ ordered by inclusion, $\varnothing$ is a subset of $\{1\}$, of $\{2\}$ and of $\{1,2\}$. So $\varnothing$ is a lower bound for $\mathcal A$. It is not an upper bound, because $\varnothing$ does not contain $\{1\}$.
\end{solution}

\begin{solution}{ex:lang-rational}
\idea The reverse implications subtract the fixed rational, or divide by it. Division needs a nonzero divisor.
\writeup First record the arithmetic. If $b,d\ne 0$, then
\[
\frac{a}{b}+\frac{c}{d}=\frac{ad+bc}{bd},\qquad
\frac{a}{b}-\frac{c}{d}=\frac{ad-bc}{bd},\qquad
\frac{a}{b}\cdot\frac{c}{d}=\frac{ac}{bd}.
\]
If also $c\ne 0$, then $(a/b)/(c/d)=ad/(bc)$. Each right-hand side is a ratio of integers with nonzero denominator. So $\Q$ is closed under addition, subtraction and multiplication, and under division by a nonzero rational.

Write $r\in\Q$. If $x\in\Q$, the sum rule gives $x+r\in\Q$. If $x+r\in\Q$, the difference rule gives $x=(x+r)-r\in\Q$. This is the first equivalence.

Now assume $r\ne 0$. If $x\in\Q$, the product rule gives $rx\in\Q$. If $rx\in\Q$, the quotient rule gives $x=(rx)/r\in\Q$, because $r\ne 0$. This is the second equivalence.

Negating both equivalences gives the irrational statements. If $x$ is irrational, then $x+r$ is irrational: a rational sum would force $x$ to be rational. If $r\ne 0$ and $x$ is irrational, then $rx$ is irrational, for the same reason read through the product. The restriction $r\ne 0$ is essential. For every real $x$ one has $0\cdot x=0$, and $0$ is rational. Multiplication by zero sends irrationals, if any, to a rational, and it also sends rationals to a rational, so it does not detect which inputs were rational.
\end{solution}

\begin{solution}{ex:lang-existence}
\idea Algebra suggests a candidate. The sign of $1-a$ checks that the candidate is positive. A second solution is forced back to the same formula.
\writeup Fix $a$ with $0<a<1$. Set $x=a/(1-a)$. The denominator $1-a$ is positive. The numerator $a$ is positive. So $x>0$, and $x$ is admissible. Also
\[
1+x=1+\frac{a}{1-a}=\frac{1-a+a}{1-a}=\frac{1}{1-a}.
\]
The denominator $1+x$ equals $1/(1-a)$, which is nonzero. Therefore
\[
\frac{x}{1+x}=\frac{a/(1-a)}{1/(1-a)}=a.
\]
This proves existence.

For uniqueness, suppose $y>0$ also satisfies $y/(1+y)=a$. The denominator $1+y$ is nonzero because $y>0$. Multiply both sides by it: $y=a(1+y)=a+ay$. Subtract $ay$: $y-ay=a$, so $(1-a)y=a$. Since $1-a>0$, we may divide by it and get $y=a/(1-a)=x$.
\end{solution}

\begin{solution}{ex:lang-falseproof}
\idea Cancellation is division, even when no division sign is written.
\writeup Let $a=b=1$. The step $a^2=ab$ is $1=1$. Subtracting $b^2$ from both sides gives $a^2-b^2=ab-b^2$, which is $0=0$. Factoring gives $(a-b)(a+b)=b(a-b)$, which is $0\cdot 2=1\cdot 0$, still true. The next step cancels the factor $a-b$. Under the opening assumption that factor equals $0$. Cancelling it divides by zero, which the arithmetic of the real numbers does not allow.

The true equation $0=0$ does not constrain $a+b$. Information discarded by a correct line cannot be recovered by an illegal cancellation. Every later numerical claim in the purported proof uses that cancellation.
\end{solution}

\begin{solution}{ex:lang-midpoint}
\idea A point outside $[a,b]$ is strictly closer to the nearer endpoint. Inside the interval the equation is linear, and it has one solution.
\writeup Let $a<b$. We first argue by cases.

If $x<a$, then $|x-a|=a-x$ and $|x-b|=b-x$. Their difference is
\[
(b-x)-(a-x)=b-a>0,
\]
so $|x-b|>|x-a|$. Equality fails.

If $x>b$, then $|x-a|=x-a$ and $|x-b|=x-b$, and $(x-a)-(x-b)=b-a>0$. Equality fails again.

If $a\le x\le b$, then $|x-a|=x-a$ and $|x-b|=b-x$. Equality becomes $x-a=b-x$, so $2x=a+b$, so $x=(a+b)/2$. This value satisfies $a<x<b$, because $a<b$ forces $(a+b)/2$ to lie strictly between them. It is the only solution of the linear equation. The three cases cover every real $x$, and only one of them contributes a solution. That proves existence and uniqueness, and identifies the point.

For the second proof, both absolute values are nonnegative. For nonnegative real numbers $u$ and $v$, one has $u=v$ if and only if $u^2=v^2$. Indeed, equal numbers have equal squares. Conversely, if $u^2=v^2$, then
\[
u^2-v^2=(u-v)(u+v)=0,
\]
so $u-v=0$ or $u+v=0$. If $u+v=0$ with $u,v\ge 0$, then $u=v=0$. If $u-v=0$, then $u=v$. In both sub-cases $u=v$.

Apply this to $u=|x-a|$ and $v=|x-b|$. Equality of the absolute values is equivalent to equality of the squares, hence to
\begin{align*}
0&=(x-a)^2-(x-b)^2\\
&=\bigl((x-a)-(x-b)\bigr)\bigl((x-a)+(x-b)\bigr)\\
&=(b-a)(2x-a-b).
\end{align*}
Since $b-a>0$, the product is zero exactly when $2x-a-b=0$, that is, when $x=(a+b)/2$. Squaring is reversible here because the two numbers being squared are nonnegative, and the factorisation above recovers equality from equality of squares.
\end{solution}

\begin{solution}{ex:lang-retraction}
\idea The equation $f(f(x))=f(x)$ says that after one application, a second application changes nothing.
\writeup (a) Suppose $y\in f(E)$. Write $y=f(x)$ for some $x\in E$. Then $f(y)=f(f(x))=f(x)=y$, so $y$ is a fixed point. Conversely, if $f(y)=y$, then $y$ is the image of itself, so $y\in f(E)$. The two sets are equal.

(b) Take $x\in E$. The element $f(x)$ is a fixed point, by (a), and $f(f(x))=f(x)$, so $f(x)$ lies in the class of $x$. Suppose $y$ and $z$ are fixed points in that same class. Then $f(y)=f(z)$, because they lie in one class, and $y=f(y)$, $z=f(z)$, so $y=z$. The fixed point in the class is unique.

(c) If $f$ is injective, the equality $f(f(x))=f(x)$ gives $f(x)=x$ for every $x$, after cancelling the outer application of $f$. If $f$ is surjective, then $f(E)=E$, so (a) says that every element of $E$ is a fixed point. In both situations $f=\operatorname{id}_E$.

For an example, take $E=\{0,1\}$ and $f(0)=f(1)=0$. Then $f(f(0))=f(0)=0$ and $f(f(1))=f(0)=0=f(1)$, so $f\circ f=f$. The function is not injective, because both inputs have the same output. It is not surjective onto $E$, because $1$ is missed. It differs from the identity, which sends $1$ to $1$. If $E=\varnothing$, the only function $E\to E$ is the empty function, which is the identity, and the conclusions hold.
\end{solution}

\section*{Challenges and onward connections}
\begin{challenge}
\textbf{1. Recover a partition from information.} Given a partition of a set $E$, map each element to its block. Prove that this map is surjective onto the set of blocks, and that its equal-output relation recovers the original partition.
\par\smallskip
\textbf{2. An order on subsets.} Prove that, under inclusion, the intersection of two subsets is their greatest lower bound, and their union is their least upper bound.
\par\smallskip
\textbf{3. A larger infinity.} Can a function $f:E\to\mathcal P(E)$ be surjective? Investigate the subset $D=\{x\in E:x\notin f(x)\}$.
\end{challenge}

\noindent\textbf{Notes.} In the first challenge, the block that contains a given element is unique, so the map is well defined. Every block is nonempty, so every block is hit. Two elements have the same image exactly when they lie in the same block, which is the original relation. Theorem~\ref{thm:lang-partition} is the same correspondence, read from the function back to the blocks.

In the second challenge, $A\cap B$ is a subset of $A$ and of $B$. If $C$ is any common subset, then every element of $C$ lies in $A$ and in $B$, so $C\subseteq A\cap B$. Thus the intersection is the greatest lower bound. The union $A\cup B$ contains $A$ and $B$. If $C$ contains both, then every element of the union lies in $C$, so $A\cup B\subseteq C$. Thus the union is the least upper bound.

For the third challenge, suppose $f(a)=D$ for some $a\in E$. By the definition of $D$, the statement $a\in D$ is equivalent to $a\notin f(a)$. Since $f(a)=D$, that statement is equivalent to $a\notin D$. The two demands $a\in D$ and $a\notin D$ cannot hold together. So no element is sent to $D$. The set $D$ is a subset of $E$, hence an element of $\mathcal P(E)$, and it is missed. The argument does not need $E$ to be nonempty. If $E=\varnothing$, then $D=\varnothing$ and $\mathcal P(E)=\{\varnothing\}$, while a function from the empty set has empty image, so it misses $D$.

This is Cantor's diagonal argument: no function from a set to its power set is surjective. Cantor gave it for the equivalent set of two-valued functions on a given set; the subset $D$ above is the same idea in modern notation. The lecture was given in September 1891. The printed paper is G.~Cantor, \emph{Ueber eine elementare Frage der Mannigfaltigkeitslehre}, \emph{Jahresbericht der Deutschen Mathematiker-Vereinigung} 1 (1890--91), 75--78; that issue appeared in 1892.

\portrait{cantor}{Georg Cantor (1845--1918). Photograph, early twentieth century. Public domain, via Wikimedia Commons.}

Volume~7, \emph{Spaces: Topology, Measure and Limits}, builds the comparison of infinite sets. The paragraph above proves only the one impossibility. It does not say that every two infinite sets can be compared by counting.

\section*{Further reading}
For further practice with the language and the organisation of proofs, see Daniel J.~Velleman, \emph{How to Prove It: A Structured Approach}, 3rd edition, Cambridge University Press, 2019, and Richard Hammack, \emph{Book of Proof}, 3rd edition, 2018, which is free to read on the author's website. These books are optional practice. They are not prerequisites.
